%% Section 5: N\'{e}ron-Tate height and height on the base

The goal of this section is to prove the full version of
Theorem~\ref{ThmHtInequality}, \textit{i.e.}, without assuming {\tt
(Hyp)}. Let $S$ be an irreducible quasi-projective variety defined over $\IQbar$ and let $\pi\colon\cA\rightarrow S$ be an abelian scheme of relative dimension $g \ge 1$. 


Let $\cL$ be a relative ample line bundle on $\cA \rightarrow S$ with $[-1]^*\cL=\cL$, and
let $\cM$ be an ample line bundle on a compactification $\overline{S}$
of $S$. All data above are assumed to be defined over $\IQbar$. Set
$\hat{h}_{\cA,\cL} \colon \cA(\IQbar) \rightarrow \IR$ to be the
fiberwise N\'{e}ron--Tate height $\hat{h}_{\cA,\cL}(P) =
\hat{h}_{\cA_s,\cL_s}(P)$ with $s = \pi(P)$, and $h_{\overline{S},\cM}
\colon \overline{S}(\IQbar) \rightarrow \IR$ to be a representative of
 the height provided by the Height Machine; \textit{cf.} \cite[Chapter~2 and 9]{BG}.

The main result of this appendix is the following theorem.

\begin{theorem}\label{ThmHtInequalityFull}
Let $X$ be an irreducible subvariety of $\cA$
defined over $\IQbar$.   Suppose $X$ is non-degenerate,
as defined in Definition~\ref{DefinitionNonDegenerateApp}. Then there
exist constants $c_1>0$ and $c_2\ge 0$ and a Zariski open dense subset
$U$ of $X$ with %and a constant $c'$ such that
\begin{equation}\label{EqHtInequalityFull}
\hat{h}_{\cA,\cL}(P) \ge c_1 h_{\overline{S},\cM}( \pi(P) ) - c_2 \quad \text{for all}\quad P \in U(\IQbar).
\end{equation}
\end{theorem}

Compared to Theorem~\ref{ThmHtInequality}, $\cA \rightarrow S$ is no longer required to satisfy {\tt (Hyp)}. Other minor improvements are that $S$ is not required to be regular and $X$ is not required to be closed.

In Definition~\ref{DefinitionNonDegenerate}, we defined non-degenerate subvarieties using the generic rank of the Betti map if $\cA \rightarrow S$ satisfies {\tt (Hyp)}. For an arbitrary $\cA \rightarrow S$, the definition is similar. But we need to first of all extend our construction of the Betti map, Proposition~\ref{PropBettiMap}, to an arbitrary $\cA \rightarrow S$.




\subsection{Betti map}
In this subsection, we extend the construction of the Betti map, \textit{i.e.}, Proposition~\ref{PropBettiMap}, to an arbitrary $\cA \rightarrow S$ with $S$ regular. %, \textit{i.e.}, without the extra assumption {\tt (Hyp)}.
In this subsection, let $S$ be an  irreducible, regular, quasi-projective variety over $\IC$ and let $\pi \colon \cA \rightarrow S$ be an abelian scheme of relative dimension $g \ge 1$.


\begin{proposition}\label{PropBettiMapApp}
Let $s_0 \in S(\IC)$. Then there exist an open neighborhood $\Delta$ of $s_0$ in $S^{\mathrm{an}}$, and a map $b_{\Delta} \colon \cA_{\Delta} := \pi^{-1}(\Delta) \rightarrow \mathbb T^{2g}$, called the Betti map, with the following properties.
\begin{enumerate}
\item[(i)] For each $s \in \Delta$ the restriction $b_{\Delta}|_{\cA_s(\IC)} \colon \cA_s(\IC) \rightarrow \mathbb T^{2g}$ is a group isomorphism.
\item[(ii)] For each $\xi \in \mathbb T^{2g}$ the preimage $b_{\Delta}^{-1}(\xi)$ is a complex analytic subset of $\cA_{\Delta}$.
\item[(iii)] The product
$(b_{\Delta},\pi) \colon \cA_{\Delta} \rightarrow \mathbb
T^{2g} \times \Delta$ is a real analytic isomorphism.% Here $\widetilde{\pi} \colon \cA_{\widetilde{S}} \rightarrow \widetilde{S}$ is the natural projection.
\end{enumerate}
\end{proposition}
%\Pinlinecomment{TODO: add a uniqueness statement along the lines of:
%  ``

Just as in the case of
Proposition~\ref{PropBettiMap}, the Betti map
is uniquely determined  by properties (i) and (iii)
up-to the action of $\mathrm{GL}_{2g}(\IZ)$ if $\Delta$ is
connected. 
% In this paper we only need the rank of the Betti map restricted to
% some subvariety.\Pcomment{I propose to strike the sentence ``In this
 %  paper\ldots'', it doesn't  seem to fit, and merge these two paragraphs.}
 Composing with an $\alpha \in \mathrm{GL}_{2g}(\IZ)$ does not change the rank. So by the discussion on the uniqueness above, any map $\cA_{\Delta} \rightarrow \mathbb{T}^{2g}$ satisfying the three properties listed in Proposition~\ref{PropBettiMapApp} will be called \textit{Betti map}.%''  I
%  think this doesn't need a proof. 
%  }
\begin{proof} Our proof of Proposition~\ref{PropBettiMapApp} follows the construction in \cite[$\mathsection$3-$\mathsection$4]{GaoBettiRank}. We divide it into several steps.

By \cite[$\mathsection$2.1]{GenestierNgo}, $\cA \rightarrow S$ carries a polarization of type $D = \mathrm{diag}(d_1,\ldots,d_g)$ for some positive integers $d_1|d_2|\cdots|d_g$.

%Let $D = \mathrm{diag}(d_1,\ldots,d_g)$ for some positive integers $d_1|d_2|\cdots|d_g$.
  \textbf{Case: Moduli space with level structure.} 
Fix $\ell \ge \bestlevel$ with $(\ell,d_g) = 1$. We start by proving Proposition~\ref{PropBettiMapApp} for $S = \mathbb
A_{g,D,\ell}$, the moduli space of abelian varieties of
dimension $g$ polarized of type $D$ with level-$\ell$-structure. It is a fine moduli
space; see \cite[Theorem~2.3.1]{GenestierNgo}. Let $\pi^{\mathrm{univ}}_D \colon \mathfrak{A}_{g,D,\ell} \rightarrow \mathbb A_{g,D,\ell}$ be the universal abelian variety.

The universal covering $\mathfrak{H}_g \rightarrow \mathbb
A_{g,D,\ell}^{\mathrm{an}}$ \cite[Proposition~1.3.2]{GenestierNgo},
where $\mathfrak{H}_g$ is the Siegel upper half space, gives a family of abelian varieties
$\cA_{\mathfrak{H}_g,D} \rightarrow \mathfrak{H}_g$ %, polarized of type $D$, 
fitting into the diagram
\[
\xymatrix{
\cA_{\mathfrak{H}_g,D} := \mathfrak{A}_{g,D,\ell} \times_{\mathbb A_{g,D,\ell}}\mathfrak{H}_g \ar[r] \ar[d] & \mathfrak{A}_{g,D,\ell}^{\mathrm{an}} \ar[d]^{\pi^{\mathrm{univ}}_D} \\
\mathfrak{H}_g \ar[r] & \mathbb A_{g,D,\ell}^{\mathrm{an}}.
}
\]
The family $\cA_{\mathfrak{H}_g,D} \rightarrow \mathfrak{H}_g$ is polarized of type $D$. 
For the universal covering
$u \colon \IC^g \times \mathfrak{H}_g \rightarrow \cA_{\mathfrak{H}_g,D}$
and for each $Z\in \mathfrak{H}_g$, the kernel of
$u|_{\mathbb{C}^g \times \{Z\}}$ is $D\mathbb{Z}^g + Z \mathbb{Z}^g$.
Thus the map
$\IC^g \times \mathfrak{H}_g \rightarrow \IR^g \times \IR^g \times \mathfrak{H}_g \rightarrow \IR^{2g}$,
where the first map is the inverse of $(a,b,Z) \mapsto (Da + Z  b,Z)$
and the second map is the natural projection, descends to a real
analytic map
\[
b^{\mathrm{univ}} \colon \cA_{\mathfrak{H}_g,D} \rightarrow \mathbb T^{2g}.
\]
 Now for each $s_0 \in \mathbb A_{g,D,\ell}(\IC)$, there exists a
 contractible, relatively compact, open neighborhood $\Delta$ of $s_0$ in $\mathbb A_{g,D,\ell}^{\mathrm{an}}$ such that $\mathfrak{A}_{g,D,\ell,\Delta}:=(\pi^{\mathrm{univ}}_D)^{-1}(\Delta)$ can be identified with $\cA_{\mathfrak{H}_g,\Delta'}$ for some open subset $\Delta'$ of $\mathfrak{H}_g$. The composite $b_{\Delta} \colon \mathfrak{A}_{g,D,\ell,\Delta} \cong \cA_{\mathfrak{H}_g,D,\Delta'} \rightarrow \mathbb T^{2g}$ clearly  satisfies the three properties listed in Proposition~\ref{PropBettiMapApp}. Thus $b_{\Delta}$ is the desired Betti map in this case.% Note that for a fixed (small enough) $\Delta$, there are infinitely choices of $\Delta'$; but for $\Delta$ small enough, if $\Delta_1'$ and $\Delta_2'$ are two such choices, then $\Delta_2' = \alpha \cdot \Delta_1'$ for some $\alpha \in \mathrm{Sp}_{2g}(\IZ) \subset \mathrm{SL}_{2g}(\IZ)$. Thus we have proven Proposition~\ref{PropBettiMapApp} for $\mathfrak{A}_g \rightarrow \mathbb A_g$.

\medskip
  \textbf{Case: With level structure.} 
  Assume that $\cA \rightarrow S$ carries level-$\ell$-structure for some $\ell \ge \bestlevel$ with $(\ell,d_g) = 1$. 
As $\IA_{g,D,\ell}$ is a fine moduli space there exists a Cartesian diagram
\begin{equation*}%\label{EqEmbedAbSchIntoUnivAbVar}
\xymatrix{
\cA \ar[r]^-{\iota} \ar[d]_{\pi} \pullbackcorner & \mathfrak{A}_{g,D,\ell} \ar[d] \\
S \ar[r]^-{\iota_S} & \mathbb{A}_{g,D,\ell}.
}
\end{equation*}

Now let $s_0 \in S(\IC)$. Applying Proposition~\ref{PropBettiMapApp}
to the universal abelian variety
$\mathfrak{A}_{g,D,\ell} \rightarrow \mathbb A_{g,D,\ell}$ and
$\iota_S(s_0) \in \mathbb A_{g,D,\ell}(\IC)$, we obtain an 
open neighborhood $\Delta_0$ of $\iota_S(s_0)$ in $\mathbb
A_{g,D,\ell}^{\mathrm{an}}$ and a real analytic map
\[
b_{\Delta_0} \colon \mathfrak{A}_{g,\Delta_0} \rightarrow \mathbb T^{2g}
\]
satisfying the properties listed in Proposition~\ref{PropBettiMapApp}.

Now let $\Delta = \iota_S^{-1}(\Delta_0)$. Then $\Delta$ is an open neighborhood of $s_0$ in $S^{\mathrm{an}}$. Denote by $\cA_{\Delta} = \pi^{-1}(\Delta)$ and define
\[
b_{\Delta} = b_{\Delta_0} \circ \iota \colon \cA_{\Delta} \rightarrow \mathbb T^{2g}.
\]
Then $b_{\Delta}$ satisfies the properties listed in Proposition~\ref{PropBettiMapApp} for $\cA \rightarrow S$. Hence $b_{\Delta}$ is our desired Betti map.


\medskip
  \textbf{Case: General case.}
Let $s_0 \in S(\IC)$ and $\ell\ge\bestlevel$ be a prime with $(\ell,d_g)=1$.

Fix any irreducible component $S_0$ of the kernel $\mathrm{ker}[\ell]$ of $[\ell] \colon
\cA\rightarrow \cA$. It is Zariski open in $\mathrm{ker}[\ell]$ as $S$ is
regular, so we consider it with its natural open subscheme structure.
Then $S_0\rightarrow \mathrm{ker}[\ell]$ is both a closed and open
immersion. So $S_0\rightarrow S$, the composition with
the finite \'etale morphism $\mathrm{ker}[\ell]\rightarrow S$, is finite and \'etale.
The upshot is that the base change of $\cA\rightarrow S$ by $S_0\rightarrow S$
admits an $\ell$-torsion section. After repeating this finitely many
times we obtain a finite and \'etale morphism $\rho \colon S'
\rightarrow S$ where $S'$ is irreducible and such that $\cA':=
\cA\times_S S' \rightarrow S'$ has level-$\ell$-structure.
Note that $S'$ is regular as $S$ is regular and regularity ascends
along \'etale morphisms. Moreover,
 $\cA' \rightarrow S'$ is still polarized of type $D$.

% There exists a finite and \'etale covering $\rho \colon S' \rightarrow
% S$, with $S'$ irreducible, such that $\cA':= \cA\times_S S' \rightarrow S'$ has
% level-$\ell$-structure.
% for some $\ell \ge \bestlevel$ with $(\ell,d_g)
% = 1$.
% \Pcomment{If we demand that $S'$  is irreducible and $\rho$ is
%  finite/\'etale,  I can see how this works without additional
%  regularity conditions on  $S$. Say $S$
%  is normal, then $\mathrm{ker}[\ell]$ is normal as
%  $\mathrm{ker}[\ell]\rightarrow S$ is \'etale. The irreducible
%  components of a normal variety are open, so we can take $S'$ as
%an irreducible component of $\mathrm{ker}[\ell]$. Then $S'\rightarrow
%S$ is finite and \'etale. See Remarques 6.15.11 in EGAIV. Are we
 %really gaining anything by dropping all regularity conditions on
%$S$ in Theorem B.1?}

Let $s_0' \in \rho^{-1}(s_0)$. Applying Proposition~\ref{PropBettiMapApp} to $\cA' \rightarrow S'$ and $s_0' \in S'(\IC)$, we obtain an open neighborhood $\Delta'$ of $s_0'$ in $(S')^{\mathrm{an}}$ and a map $b_{\Delta'} \colon \cA'_{\Delta'} \rightarrow \mathbb{T}^{2g}$ satisfying the properties listed in Proposition~\ref{PropBettiMapApp}.

Let $\Delta = \rho(\Delta')$. Up to shrinking $\Delta'$, we may assume that $\rho|_{\Delta'} \colon \Delta' \rightarrow \Delta$ is a homeomorphism and that $\Delta$ is an open neighborhood of $s_0$ in $S^{\mathrm{an}}$. Thus $\cA'_{\Delta'} \cong \cA_{\Delta}$. Now define
\[
b_{\Delta} \colon \cA_{\Delta} \rightarrow \mathbb{T}^{2g}
\]
to be the composite of the inverse of $\cA'_{\Delta'} \cong \cA_{\Delta}$ and $b_{\Delta'}$. Then $b_{\Delta}$ is our desired Betti map.
\end{proof}

%The proof of Lemma~\ref{LemmaBettiRankZarOpen} applies to this more general Betti map. We have:
Here is an easy property of the generic rank of the Betti map.
\begin{lemma}\label{LemmaBettiRankZarOpenApp}
Let $b_{\Delta} \colon \cA_{\Delta} \rightarrow \mathbb T^{2g}$ be a Betti map as in Proposition~\ref{PropBettiMapApp}. 
Let $X$ be an irreducible subvariety of $\cA$ with $\an{X}\cap\cA_\Delta\not=\emptyset$. Let $U$ be a Zariski open dense subset of $X$. Then
\begin{equation}\label{EqBettiRankSmallerUnderRestriction}
\max_{x \in X^{\mathrm{sm}}(\IC) \cap \cA_{\Delta}} \mathrm{rank}_{\IR} (\mathrm{d}b_{\Delta}|_{X^{\mathrm{sm,an}}})_x = \max_{u \in U^{\mathrm{sm}}(\IC) \cap \cA_{\Delta}} \mathrm{rank}_{\IR} (\mathrm{d}b_{\Delta}|_{U^{\mathrm{sm,an}}})_u.
\end{equation}
%\[
%\max_{\widetilde{x} \in \widetilde{X}} \mathrm{rank}_{\IR} (\mathrm{d}b|_{\widetilde{X}})_{\widetilde{x}} = \max_{\widetilde{u} \in \widetilde{U}} \mathrm{rank}_{\IR} (\mathrm{d}b|_{\widetilde{U}})_{\widetilde{u}}.
%\]
\end{lemma}
\begin{proof}
The statement \eqref{EqBettiRankSmallerUnderRestriction} is true on
replacing ``$=$'' by  ``$\ge$'', as
 $\sman{X}\supset \sman{U}$. 
%% It suffices to prove the lemma by replacing $X$ by $X^{\mathrm{sm}}$ and $U$ by $U \cap X^{\mathrm{sm}}$. Hence we may assume $X$ to be smooth. In particular $U^{\mathrm{sm}} \subseteq X$ and hence we have
%% \begin{equation}\label{EqBettiRankSmallerUnderRestriction}
%% \max_{x \in X(\IC) \cap \cA_{\Delta}} \mathrm{rank}_{\IR} (\mathrm{d}b_{\Delta}|_{X^{\mathrm{an}}})_x \ge \max_{u \in U^{\mathrm{sm}}(\IC) \cap \cA_{\Delta}} \mathrm{rank}_{\IR} (\mathrm{d}b_{\Delta}|_{U^{\mathrm{sm,an}}})_u.
%% \end{equation}

For the converse inequality we set
 $\max_{x \in \sman{X} \cap \cA_{\Delta}} \mathrm{rank}_{\IR} (\mathrm{d}b_{\Delta}|_{X^{\mathrm{sm,an}}})_x  = r$
and pick
 $x \in \sman{X} \cap \cA_{\Delta}$ satisfying
$\mathrm{rank}_{\IR} (\mathrm{d}b_{\Delta}|_{X^{\mathrm{sm,an}}})_x =
r$. Then there exists an open neighborhood  $V$ of $x$ in 
$\sman{X}$ such that $\mathrm{rank}_{\IR}
(\mathrm{d}b_{\Delta}|_{X^{\mathrm{sm,an}}})_u = r$ for all $u \in
V$. But $U^{\mathrm{sm}}(\IC) \cap V \not= \emptyset$ since
$U^{\mathrm{sm}}\not=\emptyset$ is Zariski open in $X$ and $V$ is Zariski dense
in $X$. Thus there exists a $u \in U^{\mathrm{sm}}(\IC) \cap V$. Then
we must have $\mathrm{rank}_{\IR}
(\mathrm{d}b_{\Delta}|_{U^{\mathrm{sm,an}}})_u = r$ and the lemma
follows.
%% The lemma then 
%% follows from \eqref{EqBettiRankSmallerUnderRestriction}.
%\[
%\max_{u \in U^{\mathrm{sm}}(\IC) \cap \cA_{\Delta}} \mathrm{rank}_{\IR} (\mathrm{d}b_{\Delta}|_{U^{\mathrm{sm,an}}})_u \le \max_{x \in X(\IC) \cap \cA_{\Delta}} \mathrm{rank}_{\IR} (\mathrm{d}b_{\Delta}|_{X^{\mathrm{sm,an}}})_x.\qedhere
%\]
%Now we are done.
\end{proof}


\subsection{Non-degenerate subvariety and
 Theorem~\ref{ThmHtInequality}}
 We keep the notation as in the beginning of this appendix. 
\begin{definition}\label{DefinitionNonDegenerateApp}
An irreducible subvariety $X$ of $\cA$ is said to be
\textit{non-degenerate} if there exists an open non-empty subset
$\Delta$ of $S^{\mathrm{sm},\mathrm{an}}$,
   with the Betti map
$b_{\Delta} \colon \cA_{\Delta}:=\pi^{-1}(\Delta) \rightarrow
\mathbb{T}^{2g}$ as in Proposition~\ref{PropBettiMapApp}, such that
$X^\mathrm{sm,an}\cap \cA_{\Delta}\not=\emptyset$ and 
\[
\max_{x \in X^{\mathrm{sm}}(\IC) \cap \cA_{\Delta}} \mathrm{rank}_{\IR} (\mathrm{d}b_{\Delta}|_{X^{\mathrm{sm,an}}})_x = 2\dim X.
\]%it satisfies one of the two equivalent conditions in property (iii) of Proposition~\ref{PropBettiForm}.
\end{definition}


%Let $S$ be a regular, irreducible, quasi-projective variety over $\IQbar$ and let $\pi \colon \cA \rightarrow S$ be an abelian scheme. Let $\cL$ be a relative ample line bundle on $\cA \rightarrow S$, and let $\cM$ be an ample line bundle on a compactification $\overline{S}$ of $S$. All data above are assumed to be defined over $\IQbar$. Define $\hat{h}_{\cA,\cL} \colon \cA(\IQbar) \rightarrow \IR$ to be $\hat{h}_{\cA,\cL}(P) = \hat{h}_{\cA_s,\cL_s}(P)$ with $s = \pi(P)$, and $h_{S,\cM} \colon S(\IQbar) \rightarrow \IR$ to a height function.

%\begin{theorem}\label{ThmHtInequalityFull}
%Let $X$ be an irreducible subvariety of $\cA$ defined over $\IQbar$.   Suppose $X$ is non-degenerate, as defined in Definition~\ref{DefinitionNonDegenerateApp}. Then there exist constants $c>0$ and $c'\ge 0$ and a Zariski open dense subset
%$U$ of $X$ with %and a constant $c'$ such that
%\begin{equation}\label{EqHtInequalityFull}
%\hat{h}_{\cA,\cL}(P) \ge c h_{S,\cM}( \pi(P) ) - c' \quad \text{for all}\quad P \in U(\IQbar).
%\end{equation}
%\end{theorem}

Now we are ready to prove Theorem~\ref{ThmHtInequalityFull}.

\begin{proof}[Proof of Theorem~\ref{ThmHtInequalityFull}]
  Let $\ell\ge 3$ be a prime.
  We will  reduce the current theorem to Theorem~\ref{ThmHtInequality}
  by successively assuming, in addition to the hypothesis of Theorem~\ref{ThmHtInequalityFull}, that
\begin{enumerate}
  \item[(i)] $X$ is Zariski closed in $\cA$,
  \item[(ii)] $\pi|_X \colon X\rightarrow S$ is dominant,
  \item[(iii)] $S$ is regular, 
  \item[(iv)] $\cA \rightarrow S$ is $S$-isogenous to an abelian scheme which carries a principal polarization,
  \item[(v)] $\cA \rightarrow S$ carries a principal polarization,
  \item[(vi)] $\cA$ carries a level $\ell$-structure, and
  \item[(vii)]  we have the same hypothesis as Theorem~\ref{ThmHtInequality}.% $\cA\rightarrow S$ furthermore is equipped with a closed
    % immersion $\cA \rightarrow \IP^n_{\IQbar} \times S$ over $S$
    % as in the beginning of $\mathsection$\ref{SectionSettingUpHtIneq}
    % and
    % a closed immersion
    % $\overline S\rightarrow \IP^m_{\IQbar}$.
%
    % \item[(iv)] $X$ is closed.
  % \item[(i)] $S$ is regular, 
  % \item[(ii)] $\cA \rightarrow S$ satisfies {\tt (Hyp)} (carries a principal polarization and has
  %   level-$\ell$-structure for some $\ell \ge \bestlevel$),
  % \item[(iii)] $\cA\rightarrow S$ furthermore is equipped with a closed immersion $\cA \rightarrow \IP^n_{\IQbar} \times S$ as in the beginning of $\mathsection$\ref{SectionSettingUpHtIneq},
  % \item[(iv)] $X$ is closed.
\end{enumerate}
We will proceed the proof with six d\'{e}vissage steps.
In d\'evissage step $n$ we will deduce
the theorem under the hypotheses (i),\ldots,($n-1$) 
from the theorem under the hypotheses (i),\ldots,($n$). 

% In each d\'{e}vissage step, all assumptions from the previous d\'{e}vissage steps are assumed to be satisfied.

\ul{First d\'{e}vissage: reduction to the case where $X$ is Zariski
  closed in $\cA$.}

Let $\overline{X}$ denote the
Zariski closure of $X$ in $\cA$. Then $X$ 
is a Zariski open dense subset of $\overline{X}$ and $\dim X = \dim
\overline{X}$. % Therefore $X$ is non-degenerate if and only if
% $\overline{X}$ is non-degenerate by
% Lemma~\ref{LemmaBettiRankZarOpenApp}.
Therefore, $\overline X$ is non-degenerate if $X$ is non-degenerate.
Now if
\eqref{EqHtInequalityFull} holds true on a Zariski open dense subset
$U$ of $\overline{X}$, then \eqref{EqHtInequalityFull} clearly holds
true on $U \cap X$, which is Zariski open and dense in $X$. Thus it suffices to prove \eqref{EqHtInequalityFull} with $X$ replaced by $\overline{X}$. %So we may assume that (iv) is furthermore satisfied.


\ul{Second d\'{e}vissage: reduction to the case where $\pi|_{X}\colon X\rightarrow 
  S$ is dominant.} 
  
  As $X$ is non-degenerate, there exists a non-empty open subset $\Delta$ of $S^{\mathrm{sm,an}}$, with Betti map $b_{\Delta}$, such that $\mathrm{rank}_{\IR}(\mathrm{d}b_{\Delta}|_{X^{\mathrm{sm,an}}})_x = 2 \dim X$ for some $x \in X^{\mathrm{sm}}(\IC) \cap \cA_{\Delta}$.

Endow the Zariski closed set $S'=\pi(X)$ with the reduced induced
subscheme structure and set $\cA' = \cA \times_S S' =
\pi^{-1}(S')$. Then $X \times_S S'$ identifies with $X$ via the natural projection $\cA' \rightarrow \cA$. Hence there exists a non-empty open subset $\Delta'$
of $(S')^{\mathrm{sm,an}}$ with $\pi(x) \in \Delta'\subset \Delta$ and
\[
 \mathrm{rank}_{\IR} (\mathrm{d}b_{\Delta'}|_{X^{\mathrm{sm,an}}})_x = 2\dim X.
\]
Thus $X$ is a non-degenerate subvariety of $\cA'$. On the other hand,
the conclusion of  Theorem~\ref{ThmHtInequalityFull} does not change
with $\cA \rightarrow S$ replaced by $\cA' \rightarrow S'$, $\cL$
replaced by $\cL|_{\cA'}$ and $\cM$ replaced by $\cM|_{\overline{S'}}$,
 where $\overline{S'}$ is the Zariski closure of $S'$ in $\overline{S}$.
Hence it suffices to prove Theorem~\ref{ThmHtInequalityFull} after
these replacements and thus we may assume that $X$ dominates $S$.


\ul{Third d\'{e}vissage: reduction to the case where $S$ is regular.}

Recall that $S^{\mathrm{sm}}$ is the regular locus of $S$.
Now $\pi|_X\colon X\rightarrow S$ is dominant, so
$X'=X\cap \pi^{-1}(S^{\mathrm{sm}})$ is Zariski open and dense in
$X$. 
Since $X$ is non-degenerate it follows by definition that  $X'$ is
non-degenerate. Moreover, the conclusion of
Theorem~\ref{ThmHtInequalityFull} does not change if we replace $\cA
\rightarrow S$ by $\cA' = \pi^{-1}(S^{\mathrm{sm}})
\rightarrow S^{\mathrm{sm}}$, $\cL$ by $\cL|_{\cA'}$,
and $X$ by $X'$.
Finally, observe that $X'$ is Zariski closed in $\cA'$
and $\pi(X') = \pi(X)\cap S^{\mathrm{reg}}$, so $\pi|_{X'}\colon X'
\rightarrow S^{\mathrm{reg}}$ is dominant. 
%So we may assume that (i) is satisfied.


\ul{Fourth d\'{e}vissage: reduction to the case where $\pi \colon \cA \rightarrow S$ is $S$-isogenous to an abelian scheme which carries a principal polarization.}

%satisfies {\tt (Hyp')}, with}
%\begin{center}
%{\tt (Hyp'):} 
%it has level-$\ell$-structure for some $\ell \ge \bestlevel$, and for its generic fiber $A$ there exists an isogenous $A_0 \rightarrow A$ defined over $\IQbar(S)$ with $A_0$ principally polarized.
%\end{center}

%There exists a quasi-finite \'{e}tale dominant %\Pcomment{Replace finite by quasi-finite?}
%morphism $S' \rightarrow S$ such that
%$\cA':= \cA\times_S S' \rightarrow S'$ has level-$\ell$-structure for
%some $\ell \ge \bestlevel$. 
By \cite[$\mathsection$23,
Corollary~1]{MumfordAbVar}, each abelian variety over an algebraic
closed field is isogenous to a principally polarized one. Applying
this to the geometric generic fiber of $\cA \rightarrow S$, we obtain a
quasi-finite \'{e}tale dominant morphism 
 $\rho \colon S' \rightarrow S$ with $S'$ irreducible and the
 following property: There exists a principally polarized $A'_0$ that
 is isogenous over $\IQbar(S')$ to  the
 generic fiber $A'$ of $\cA':= \cA\times_{S} S' \rightarrow S'$. 
 Up to replacing $S'$ by an open dense subscheme, we may furthermore assume that $A'_0$ extends to an abelian scheme $\cA'_0 \rightarrow S'$. 
 Denote by $\rho_{\cA} \colon \cA' =\cA \times_S S' \rightarrow \cA$ the natural projection; it is a quasi-finite \'{e}tale dominant morphism.

As regularity ascends along \'etale morphisms and as $S$ is regular we conclude that $S'$ is regular. Thus $\cA'_0 \rightarrow S'$ carries a principal polarization by \cite[Th\'{e}or\`{e}me~XI~1.4]{LNM119}, and the isogeny $A'_0 \rightarrow A'$ extends to an $S'$-isogeny $\cA'_0 \rightarrow \cA'$ by \cite[Lemme~XI~1.15]{LNM119}.% Thus  the isogeny from $A'_0$ to the generic fiber of $\cA' \rightarrow S'$ extends to an $S'$-isogeny $\cA'_0 \rightarrow \cA'$, \textit{cf.} \cite[Th\'{e}or\`{e}me~XI~1.4]{LNM119}.% The morphism $\lambda$ is a finite morphism, and hence $\lambda \circ \rho_{\cA} \colon \cA'_0 \rightarrow \cA$ is quasi-finite dominant.


There is  an irreducible component $X'$ of $\rho_{\cA}^{-1}(X)$ with
$\dim X' = \dim X$. Then $X'$ is Zariski closed in $\cA'$, 
the image $\rho_{\cA}(X')$ is Zariski dense in $X$,  and thus $X'$ dominates $S'$ (it even surjects to $S'$
since $\cA'\rightarrow S'$ is proper and $X'$ is closed). We
claim that $X'$, as a subvariety of $\cA'$, is non-degenerate.
Indeed, $\rho_{\cA}(X')$ contains a Zariski open dense subset $U$ of
$X$. Since $X$ is a non-degenerate
subvariety of $\cA$, so is $U$ by
Lemma~\ref{LemmaBettiRankZarOpenApp}. So there exists an open subset
$\Delta$ of $\an{S}$ with the Betti map $b_{\Delta} \colon
\cA_{\Delta} \rightarrow \mathbb{T}^{2g}$ such that
$\mathrm{rank}_{\IR}(\mathrm{d}b_{\Delta}|_{\ansm{U}})_u = 2 \dim U =
2 \dim X$ %\Pinlinecomment{
for all $u$ from a non-empty open subset of $\an{U}$
%}
. Take $\Delta'$ to be a connected component of $\rho^{-1}(\Delta)$ such that $X'
\cap (\pi')^{-1}(\Delta') \not= \emptyset$. Set $\cA'_{\Delta'} =
(\pi')^{-1}(\Delta')$, and replace $\Delta$ by $\rho(\Delta')$.
Note that $\rho|_{\Delta'}\colon \Delta' \cong \Delta$ is then
bianalytic after possibly shrinking $\Delta'$ (and so is $\rho_{\cA}\colon \cA'_{\Delta'} \cong
\cA_{\Delta}$). Now $b_{\Delta} \circ \rho_{\cA}|_{\cA'_{\Delta'}} \colon \cA'_{\Delta'} \rightarrow \mathbb{T}^{2g}$ satisfies the three properties listed in Proposition~\ref{PropBettiMapApp}. So $b_{\Delta} \circ \rho_{\cA}|_{\cA'_{\Delta'}}$ is the Betti map, which we denote for simplicity by $b_{\Delta'}$; see below Proposition~\ref{PropBettiMapApp}. 
For $u' \in  (\rho_{\cA}|_{\cA'_{\Delta'}})^{-1}(u)\cap \an{X'}$
 and for sufficiently general $u$%\Pcomment{Slight
%  elaboration as $\rho|_{X'}$ must have maximal rank at $u$}
, we have 
$\mathrm{rank}_{\IR}(\mathrm{d}b_{\Delta'}|_{\ansm{X'}})_{u'} = 2 \dim X$. So $X'$, as a subvariety of the abelian scheme $\cA'$ over $S'$, is non-degenerate.

Now we have a non-degenerate subvariety $X'$ of the abelian scheme $\pi' \colon \cA' \rightarrow S'$. The line bundle $\rho_{\cA}^*\cL$ on $\cA'$ is relatively ample. Suppose that $\cM'$ is an ample line bundle on some compactification $\overline{S'}$ of $S'$.

Assume that Theorem~\ref{ThmHtInequalityFull} holds for $\pi' \colon \cA' \rightarrow S'$, $\rho_{\cA}^*\cL$, $\cM'$, and $X'$. Thus there exist constants $c_1'>0$, $c_2' \ge 0$ and a Zariski open non-empty subset $U'$ of $X'$ with
\[
\hat{h}_{\cA',\rho_{\cA}^*\cL}(P') \ge c_1' h_{\overline{S'},\cM'}(\pi'(P')) - c_2' \quad \text{for all} \quad P' \in U'(\IQbar).
\]
Denote by $P = \rho_{\cA}(P')$. 
By the Height Machine we have $\hat{h}_{\cA',\rho_{\cA}^*\cL}(P') =
\hat{h}_{\cA,\cL}(P)$.%\Pinlinecomment{We have to be more careful when
%  applying the height machine. Our reference
%  Bombieri-Gubler doesn't define the  HM when the variety in
%question is only quasi-projective. Projectiveness ist crucial for the
%functoriality properties.} \Gcomment{The $\hat{h}$ is fiberwise
%defined in our convention. As fiberwise the variety is projective, I
%think it's ok to apply the Height Machine
%directly.}\Pinlinecomment{This argument may work for $\hat h$, but not
%for $h_{S'',\cM''}$ for which  we freely apply HM arguments below. I
%think it's easy to fix this, will look into it later.} \Ginlinecomment{I don't think we apply HM to $h_{S'',\cM''}$ below. I changed the notation on height function $h_{S,\cM}$ to $h_{\overline{S},\cM}$ to emphasize that $\cM$ is an ample line bundle on $\overline{S}$ and not just on $S$.}

By  \cite[Lemma~4]{Silverman:heightest11} applied to $\rho \colon S' \rightarrow S$ and the line bundles $\cM'$ and $\cM$, there exist $c' = c' (\rho,\cM',\cM) > 0$ and $c'' = c''(\rho,\cM',\cM) \ge 0$ such that $h_{\overline{S'},\cM'}(\pi'(P')) \ge c' h_{\overline{S},\cM}(\rho(\pi'(P'))) - c'' = c' h_{\overline{S},\cM}(\pi(P)) - c''$ for all $P' \in \cA'(\IQbar)$. Hence the height inequality above implies
\[
\hat{h}_{\cA,\cL}(P) \ge c_1'c' h_{\overline{S},\cM}(\pi(P)) - (c_1'c'' + c_2')\quad \text{for all}\quad P \in \rho_{\cA}(U')(\IQbar).
\]
Now that $\rho_{\cA}(U')$ contains a Zariski open non-empty (hence dense) subset $U$ of $X$ by Chevalley's Theorem. 
Thus Theorem~\ref{ThmHtInequalityFull} also holds true for $\pi \colon \cA \rightarrow S$, $\cL$, $\cM$, and $X$.

In summary, we have shown that it suffices to prove Theorem~\ref{ThmHtInequalityFull} for $\pi' \colon \cA' \rightarrow S'$, $\rho_{\cA}^*\cL$, $\cM'$ and $X'$. Thus we are reduced to the case where the generic fiber of $\cA \rightarrow S$ is isogenous to a principally polarized abelian variety.%carries a principal polarization.


\ul{Fifth d\'{e}vissage: reduction to the case where $\pi \colon \cA \rightarrow S$ carries a principal polarization.}

From the previous d\'{e}vissage, there exists a principally polarized abelian scheme $\pi_0 \colon \cA_0 \rightarrow S$ with an $S$-isogeny $\lambda \colon \cA_0 \rightarrow \cA$. Note that $\lambda$ is a finite \'{e}tale morphism. The line bundle $\lambda^*\cL$ on $\cA_0$ is relatively ample. By the Height Machine we have $\hat{h}_{\cA_0,\lambda^*\cL}(P') = \hat{h}_{\cA,\cL}(\lambda(P'))$ for all $P' \in \cA_0(\IQbar)$.

There is  an irreducible component $X_0$ of $\lambda^{-1}(X)$ with
$\dim X_0 = \dim X$. Then $X_0$ is Zariski closed in $\cA_0$ and thus $X_0$ dominates $S$ (it even surjects to $S$
since $X = \lambda(X_0)$). We
claim that $X_0$, as a subvariety of $\cA_0$, is non-degenerate. Assume this. Then it suffices to prove the height inequality \eqref{EqHtInequalityFull} with $\cA \rightarrow S$ replaced by $\cA_0 \rightarrow S$, $X$ replaced by $X_0$, and $\cL$ replaced by $\lambda^*\cL$. 

It remains to prove that $X_0$ is a non-degenerate subvariety of $\cA_0$. To do this, we need some preparation on Betti maps. Let $\Delta$ be an open subset of $S^{\mathrm{an}}$ with the Betti map $b_{\Delta} \colon \cA_{\Delta} \rightarrow \mathbb{T}^{2g}$. Set $\cA_{0,\Delta} = \pi_0^{-1}(\Delta)$, and denote by $\lambda_{\Delta}$ the restriction of $\lambda \colon \cA_0 \rightarrow \cA$ to $\cA_{0,\Delta}$. Up to shrinking $\Delta$ we have a Betti map $b_{0,\Delta} \colon \cA_{0,\Delta} \rightarrow \mathbb{T}^{2g}$. By property (iii) of Proposition~\ref{PropBettiMapApp}, we have two real analytic isomorphisms $(b_{0,\Delta},\pi_0) \colon \cA_{0,\Delta} \cong \mathbb{T}^{2g} \times \Delta$ and $(b_{\Delta},\pi) \colon \cA_{\Delta} \cong \mathbb{T}^{2g} \times \Delta$. Thus there exists a real analytic map $\lambda' \colon \mathbb{T}^{2g} \times \Delta \rightarrow \mathbb{T}^{2g} \times \Delta$ such that the following diagram commutes
\[
\xymatrix{
\cA_{0,\Delta} \ar[r]^-{(b_{0,\Delta},\pi_0)}_-{\sim} \ar[d]_-{\lambda_{\Delta}} & \mathbb{T}^{2g} \times \Delta \ar[d]^-{\lambda'} \\
\cA_{\Delta} \ar[r]^-{(b_{\Delta},\pi)}_-{\sim} & \mathbb{T}^{2g} \times \Delta.
}
\]
As $\lambda$ is a finite map, $(\lambda')^{-1}(r)$ is a finite set for each $r \in \mathbb{T}^{2g} \times\Delta$. As $\lambda$ is an $S$-morphism, for each $s \in \Delta$ we have $\lambda'({\mathbb{T}^{2g} \times \{s\}}) \subseteq \mathbb{T}^{2g} \times \{s\}$.

By property (i) of Proposition~\ref{PropBettiMapApp}, for each $s \in \Delta$ the restriction $\lambda'|_{\mathbb{T}^{2g} \times \{s\}}$ is a group homomorphism $\mathbb{T}^{2g}  \rightarrow  \mathbb{T}^{2g}$. Thus $\ker(\lambda'|_{\mathbb{T}^{2g} \times \{s\}})$ is a finite, hence discrete, subgroup of $\mathbb{T}^{2g}$. In particular, $\ker(\lambda'|_{\mathbb{T}^{2g} \times \{s\}})$ is locally constant. Up to shrinking $\Delta$, we may assume $\ker(\lambda'|_{\mathbb{T}^{2g} \times \{s\}}) = H$ for each $s \in \Delta$. Set $\lambda_{\mathbb{T}} \colon \mathbb{T}^{2g} \rightarrow \mathbb{T}^{2g}$ the quotient by the finite subgroup $H$. 
Then the diagram above induces a commutative diagram
\begin{equation}
\xymatrix{
\cA_{0,\Delta} \ar[r]^-{b_{0,\Delta}} \ar[d]_-{\lambda_{\Delta}} & \mathbb{T}^{2g} \ar[d]^-{\lambda_{\mathbb{T}}} \\
\cA_{\Delta} \ar[r]^-{b_{\Delta}} & \mathbb{T}^{2g}.
}
\end{equation}
Note that $\lambda_{\mathbb{T}}$ is a local homeomorphism.


Now we turn to proving that $X_0$ is non-degenerate in $\cA_0$. Indeed, as $X$ is a non-degenerate subvariety of $\cA$, there exists an open subset $\Delta$ of $S^{\mathrm{an}}$ with the Betti map $b_{\Delta} \colon \cA_{\Delta} \rightarrow \mathbb{T}^{2g}$ such that $\mathrm{rank}_{\IR}(\mathrm{d}b_{\Delta}|_{X^{\mathrm{an,sm}}})_x = 2 \dim X$ for all $x$ from a non-empty open subset of $X^{\mathrm{an}}$. For $x_0 \in  \lambda^{-1}(x)\cap X_0^{\mathrm{an}}$
 and for sufficiently general $x$%\Pcomment{Slight
%  elaboration as $\rho|_{X'}$ must have maximal rank at $u$}
, we have 
$\mathrm{rank}_{\IR}(\mathrm{d}(\lambda_{\mathbb{T}}\circ b_{0,\Delta})|_{X_0^{\mathrm{an,sm}}})_{x_0} = 2 \dim X = 2 \dim X_0$. But $\lambda_{\mathbb{T}}$ is a local homeomorphism, so $\mathrm{rank}_{\IR}(\mathrm{d} b_{0,\Delta}|_{X_0^{\mathrm{an,sm}}})_{x_0} = 2 \dim X_0$. Thus $X_0$ is non-degenerate.



% $A_0$ that is isogenous over $\IQbar(S)$ to the generic fiber of $\cA \rightarrow S$. As $S$ is regular, we can apply \cite[Th\'{e}or\`{e}me~XI~1.4]{}


\ul{Sixth d\'{e}vissage: reduction to the case where $\cA/S$
  carries level $\ell$ structure.}

As in the treatment of the general case in the proof of
Proposition~\ref{PropBettiMapApp}
there exists a finite and \'etale morphism $S'\rightarrow S$ where
$S'$ is regular and irreducible
 such that $\cA' := \cA\times_S S'$ carries level $\ell$-structure.
%By {\tt (Hyp')}, there exists a Zariski open dense subset $S_0$ of $S$ such that: There exists an abelian scheme $\cA_0 \rightarrow S_0$, carrying a principal polarization, with an $S_0$-isogeny $\lambda \colon \cA_0 \rightarrow \cA_{S_0} := \pi^{-1}(S_0)$. 
%Up to replacing $S_0$ by its smooth locus, we may furthermore assume
%that $S_0$ is regular. Moreover, $\cA_0 \rightarrow S_0$ has
%level-$\ell$-structure for some $\ell \ge \bestlevel$ because
%$\cA_{S_0} \rightarrow S_0$ does.\Pcomment{I tried to verify this claim
%  by pulling back an $\ell$-torsion section using the
%  $S_0$-isogeny $\lambda$. But the kernel of the isogeny could
%  conceivably contain this section, e.g. if $\ell$ divides the  degree
%  of the isogeny. In this case it
%  seems that the level structure is spoiled. Can we do the isogeny
%  first and then introduce level structure.}


%Now that $S$ is regular and $\cA \rightarrow S$ carries a principal polarization. 
%!!!!
%Then $S'$ is regular for the same reason as in the fourth d\'evissage.
Denote by $\rho_{\cA} \colon \cA' \rightarrow \cA$ the natural
projection. By a similar argument as the fourth d\'{e}vissage step, it
suffices to prove the height inequality \eqref{EqHtInequalityFull} with
$\cA \rightarrow S$ replaced by $\cA' \rightarrow S'$, $X$ replaced by
an irreducible component of $\rho_{\cA}^{-1}(X)$ with $\dim X'=\dim
X$, and $\cL$ replaced by $\rho_{\cA}^*\cL$, and $\cM$ replaced by
an ample line bundle on some compactification of $S'$.
% (notice that $\hat
  % h_{\cA',\rho_{\cA}^*\cL} = \hat
  % h_{\cA,\cL}\circ \rho_{\cA}$ by the Height Machine).
  As in the fourth d\'evissage  $X'$ dominates $S'$. Finally,
  $\cA'/S'$ still carries a principal polarization.
  %So we may assume that (ii) is furthermore satisfied. 
%\Pinlinecomment{I guess we also have to bound $\hat  h_{\cA_0,\lambda^*\cL}$ from above in terms of $\hat   h_{\cA,\cL}\circ\lambda$, but the two are equal by functoriality   of the height. Maybe we can add a short comment to this extent.} 


\ul{Seventh d\'{e}vissage: reduction to  Theorem {\ref{ThmHtInequality}}.}

It remains to prove the height inequality (\ref{EqHtInequalityFull})
with the extra hypotheses (i) - (v) listed above using
Theorem~\ref{ThmHtInequality}. In this theorem we assumed in addition that 
the
fiberwise  N\'eron--Tate height on $\cA(\IQbar)$
is induced by  a closed immersion $\cA \rightarrow \IP^n_{\IQbar}
\times S$ satisfying the second and third bullet  at the beginning of
$\mathsection$\ref{SectionSettingUpHtIneq} and
that the height on $\overline S(\IQbar)$ is the restriction of the absolute
logarthmic Weil height coming from a closed immersion
$\overline S\rightarrow\IP^m_{\IQbar}$. 
% % \Pinlinecomment{There may be some
% %  overlap with Remark \ref{rmk:projembedding}}
% First of all note that the fourth bullet point already holds true.
% Indeed, \texttt{(Hyp)} were handled in the fourth and fifth d\'evissages.

%\Pcomment{I slightly reworded the next paragraph.}
A basis of the global sections of the line bundle $\cM^{\otimes p}$,
for some $p$ large enough, gives rise to a closed immersion $\overline S \subseteq
\IP_{\IQbar}^m$. This gives the first bullet point at the beginning of
$\mathsection$\ref{SectionSettingUpHtIneq}. Note that the Weil
height $h$ on $\IP_\IQbar^m(\IQbar)$ restricted to $\overline
S(\IQbar)$ via this immersion differs from 
$p h_{\overline S,\cM}$ by a bounded function.
% on $S$ given by this immersion $S \subseteq \IP^m$, we have
% $h_{\overline S,\cM} = (1/p)h$ up to a bounded function on $S(\IQbar)$.

For the line bundle $\cL$ on $\cA$, which is ample relative over $S$, we have that $\cL^{\otimes 4}$ is relatively very ample on $\cA/S$. Thus by  \cite[Proposition~4.4.10.(ii) and Proposition~4.1.4]{EGAII}, there is a closed immersion $\cA \rightarrow \IP^n_S = \IP^n \times S$ given by global sections of $\cL^{\otimes 4} \otimes \pi^*\cM^{\otimes q}$ for some large $q$. When restricted to the generic fiber $A$ of $\cA \rightarrow S$, we get a closed immersion $A \rightarrow \IP^n_{k(S)}$ which arises from a basis of the global sections of $L^{\otimes 4}$, where $L$ is the restriction of $\cL$ over the generic fiber $A$. Moreover $L$ is ample since $\cL$ is relatively ample, and $L$ is symmetric since $[-1]^*\cL = \cL$. Thus we also have the second and third bullet points at the beginning of $\mathsection$\ref{SectionSettingUpHtIneq}.
%Let $\eta$ be the generic fiber of $\cA \rightarrow S$. The desired immersion can be obtained by applying Remark~\ref{rmk:projembedding} to $L_0 = \cL_{\eta}$; we point out that the $\cL$, $\cM$ in Remark~\ref{rmk:projembedding} are different from our $\cL$, $\cM$ here. 

Note that the height function $\hat{h}_{\cA}$ defined in \eqref{eq:fiberwiseNT} is then
\[
\hat{h}_{\cA} \colon \cA(\IQbar) \rightarrow [0,\infty), \quad P \rightarrow \hat{h}_{\cA_s,\cL_s^{\otimes 4}}(P)
\]
where $s = \pi(P)$. So $\hat{h}_{\cA,\cL} = (1/4)\hat{h}_{\cA}$. 


The full
hypothesis of Theorem~\ref{ThmHtInequality}  is now satisfied for $\cA$
and $X$,
\textit{e.g.}, \texttt{(Hyp)} is just (iv) and (v).
We get constants $c_1> 0$ and $c_2$ and a Zariski open dense subset $U$ of $X$ such that
\[
    \hat{h}_{\cA}(P) \ge c_1 h( \pi(P) ) - c_2 \quad \text{for all}\quad P \in U(\IQbar).
\]
Thus \eqref{EqHtInequalityFull} holds true with $c_1$ replaced by $(c_1p)/4$ and $c_2$ replaced by $c_2/4 + O_S(1)$, where $O_S(1)$ is a bounded function on $S(\IQbar)$. So we are done.
\end{proof}



%%% Local Variables:
%%% TeX-master: "main"
%%% End:
